\scanpage{83}
% archive foliation: 82

de $\kappa$). On peut dire aussi que la donnée de $\kappa$ équivaut à celle d'une \deleted{paire}\inserted{couple} $(\alpha,\kappa)$, où $\alpha$ est une section de $\unsure{\boldsymbol{\mu}|\mu}_4^*\wedge\underline{\omega}_{\underline{J}}$, et où $\kappa$ est un iso. (31) \uline{compatible} avec $\alpha$ au sens précédent. Pour $\alpha$ fixé, les $\kappa$ correspondants forment les sections d'un torseur sous $\unsure{\boldsymbol{\mu}|\mu}_2$ --- un \uline{torseur des épinglages de Legendre} de $E/S$, \deleted{\unsure{au sens de}}\inserted{associé à} ``l'épinglage partiel'' $\alpha$.

La catégorie\deleted{groupoïde}\inserted{fibrée} (qui est un champ en groupoïdes) des courbes elliptiques relatives, munies d'un épinglage de Legendre, équivaut à celle des triplets $(\underline{J},t,\alpha)$, où
\[
(34)\quad\left\{\begin{array}{ll}
a)&\underline{J}\text{ rev. étale cubique de }S\\[4pt]
b)&t\text{ section de }\Sigma_{\underline{J}}^*\\[4pt]
c)&\alpha\text{ section de }\unsure{\boldsymbol{\mu}|\mu}_4^*\wedge\underline{\omega}_{\underline{J}},\ \text{i.e. iso. }\unsure{\boldsymbol{\mu}|\mu}_4^*\simeq\underline{\omega}_{\underline{J}}
\end{array}\right.
\]
\margin{\uline{NB} Il faudrait expliciter la construction de $E$, en termes de ces données\unsure{\ldots} (via \unsure{fonctions} $\wp$ de Weierstrass)}
